Numbers are read in $\mathrm{span}\{1, I_{4}\}$ throughout, and $\bar{z}$ is $I_{4} \mapsto -I_{4}$.

\emph{(i).} For bivectors $X$, $Y$ the product $XY$ has grades $0$, $2$, $4$, its grade-two part is the commutator $\tfrac12(XY - YX)$ and its grade-four part is $X \wedge Y$, symmetric in $X$ and $Y$; so $\tfrac12\langle XY + YX\rangle_{4} = X \wedge Y$. The relative vectors $\gamma_{i}\gamma_{0}$ anticommute pairwise and square to $1$, $I_{4}$ commutes with them and $I_{4}^{2} = -1$, so $(E + I_{4}B)^{2} = E^{2} - B^{2} + I_{4}(EB + BE) = \sum_{i}(E_{i}^{2} - B_{i}^{2}) + 2\bigl(\sum_{i}E_{i}B_{i}\bigr)I_{4}$; hence $\langle I_{4}F^{2}\rangle_{0} = -2\sum_{i}E_{i}B_{i}$, the form of \claimref{synthesis.triangle-anomaly-in-the-spacetime-algebra-is-the-pseudoscalar-part-of-the-field-square-and-cancels-per-generation}(iv). Reversal changes the sign of a bivector and the order of a product, so it sends $\tfrac12(XY + YX)$ to itself, and it fixes grade four; the grade involution fixes even elements. An outermorphism preserves the outer product, and on the top grade it is multiplication by the determinant: $g(X) \wedge g(Y) = g(X \wedge Y) = \det(g)\,X \wedge Y$. Summing with the constants $\kappa_{ab}$ gives the three statements for $\Theta$, which is the invariant polynomial $\tr$ evaluated on the field square, the four-form of \claimref{synthesis.imported-chern-weil-and-chern-simons-forms} with $k = 2$. The check verifies them on rational bivectors for a reflection, the time reversal, the total inversion, a boost and a rotation.

\emph{(ii).} By clause (1) of \claimref{synthesis.cpt-conjugation-pseudoscalar-rotor-times-complex-conjugation} the grade involution equals $\mathrm{Ad}(I_{6})$, and by clause (4) $C = \mathrm{Ad}(I_{6}) \circ \kappa$ is $-1$ on $\R^{2,4}$ and $\kappa$ fixes $\Cl(2,4)$. The map $-1$ of a six-dimensional space has determinant $+1$, and its outermorphism is $(-1)^{k}$ on grade $k$: the identity on bivectors and on four-vectors. In the convention of clause (5) a field built from vectors is acted on through the vector part, so every bivector field and every four-form built from two of them is fixed. On the colour index $C$ acts, in the same convention, through the Lie algebra automorphism it induces. If $\tau$ is an automorphism of $\su(3)$ then $\mathrm{ad}(\tau X) = \tau\,\mathrm{ad}(X)\,\tau^{-1}$, so $\tr(\mathrm{ad}\,\tau X\,\mathrm{ad}\,\tau Y) = \tr(\mathrm{ad}\,X\,\mathrm{ad}\,Y)$; the trace form of the defining representation is a multiple of this form, $\su(3)$ being simple, so it is preserved, whichever automorphism the conjugation induces (the check verifies it for the outer one, $X \mapsto \bar{X}$, and for an inner one). The trace form is real on $\su(3)$, so the antilinearity of $\kappa$ does not act on $\kappa_{ab}$. Hence $\Theta(CF) = \Theta(F)$. A real-linear flag conjugate $g$ has $\det g = -1$ by clause (3) there, so $g(I_{6}) = -I_{6}$; if it preserves a four-dimensional subspace and has determinant $-1$ on it, (i) applied in that subspace gives $\Theta(gF) = -\Theta(F)$ for $F$ in it. The check verifies on $\Cl(2,4)$ that $\mathrm{Ad}(I_{6})$ is the grade involution, that a reflection and the product of three reflections reverse $I_{6}$, and that a reflection reverses exactly the four-forms containing its direction.

\emph{(iii).} The mirror tower is related to this one by the signature reversal of \claimref{dark-matter.mirror-tower-copy-one}. Put $b_{\mu\nu} = \gamma_{\mu}\gamma_{\nu}$ and $b'_{\mu\nu} = \gamma'_{\mu}\gamma'_{\nu}$ for $\mu \neq \nu$. The even subalgebra is generated by the $b_{\mu\nu}$ with the relations $b_{\mu\nu}b_{\nu\rho} = Q(\gamma_{\nu})\,b_{\mu\rho}$ for distinct indices, $b_{\mu\nu}^{2} = -Q(\gamma_{\mu})Q(\gamma_{\nu})$, $b_{\nu\mu} = -b_{\mu\nu}$, and the commutation of bivectors on disjoint index pairs. Under $b_{\mu\nu} \mapsto -b'_{\mu\nu}$ the first relation becomes $b'_{\mu\nu}b'_{\nu\rho} = -Q(\gamma_{\nu})b'_{\mu\rho} = Q(\gamma'_{\nu})b'_{\mu\rho}$, which holds; the second holds because both squares change sign twice; the others are linear. So the map extends to an algebra isomorphism $\sigma$, and $\sigma(I_{4}) = \sigma(b_{01}b_{23}) = b'_{01}b'_{23} = I'_{4}$. An algebra isomorphism that maps bivectors to bivectors and $I_{4}$ to $I'_{4}$ maps the grade-four part of a product of two bivectors to the grade-four part of the product of the images. The map $b_{\mu\nu} \mapsto b'_{\mu\nu}$ fails the first relation; it preserves outer products of basis blades, and since $\Theta$ is quadratic the sign on bivectors does not enter. The check verifies multiplicativity on the whole basis of the even algebra and the equality of the two four-forms.

\emph{(iv).} The gauge group acts on $\Phi_{1} = \alpha$ and $\Phi_{2} = \bar\beta$ by the same doublet action, hypercharge by a common phase ($\alpha \mapsto X\alpha$, $\beta \mapsto \bar{X}\beta$ in the cited Theorem). A real-linear map of the rung summand commuting with those phases is linear over $\mathrm{span}\{1, I_{4}\}$, and one commuting with the irreducible doublet action on the two copies is a $2 \times 2$ matrix acting on the pair; it keeps $\mathsf{N} = \tfrac12(|\Phi_{1}|^{2} + |\Phi_{2}|^{2})$ exactly when it is unitary. Let $A_{1} = Y$, $A_{2} = \tilde{Y}$, $B_{1} = \tilde{Y}$, $B_{2} = Y$, so $M_{u} = \sum_{j}\Phi_{j}A_{j}$ and $M_{d} = \sum_{j}\bar\Phi_{j}B_{j}$, and let $K_{\Phi} = (k_{ij})$ act by $\delta\Phi_{i} = \sum_{j}k_{ij}\Phi_{j}$. The coefficients of $\Phi_{j}$ and of $\bar\Phi_{j}$ in the two conditions give, for $j = 1, 2$,
\[
\textstyle\sum_{i}k_{ij}A_{i} + K_{Q}^{\mathsf T}A_{j} + A_{j}K_{u} = 0, \qquad \sum_{i}\bar{k}_{ij}B_{i} + K_{Q}^{\mathsf T}B_{j} + B_{j}K_{d} = 0 .
\]
Write $S = 2\tr K_{Q} + \tr K_{u} + \tr K_{d}$. Suppose $\det Y \neq 0$. Multiply the first equation at $j = 1$ and the second at $j = 2$ by $Y^{-1}$ and take traces: $nk_{11} + k_{21}\tr(Y^{-1}\tilde{Y}) + \tr K_{Q} + \tr K_{u} = 0$ and $\bar{k}_{12}\tr(Y^{-1}\tilde{Y}) + n\bar{k}_{22} + \tr K_{Q} + \tr K_{d} = 0$. $K_{\Phi}$ is anti-Hermitian: $k_{21} + \bar{k}_{12} = 0$, $k_{11} = I_{4}q_{1}$, $\bar{k}_{22} = -I_{4}q_{2}$. Adding, $S = -nI_{4}(q_{1} - q_{2})$, that is $a_{3} = -n(q_{1} - q_{2})$. If also $\det\tilde{Y} \neq 0$, the first equation at $j = 2$ and the second at $j = 1$, multiplied by $\tilde{Y}^{-1}$, give in the same way $S = -nI_{4}(q_{2} - q_{1})$. So $S = 0$ and $q_{1} = q_{2}$. At $\alpha = 0$, $M_{d} = \beta Y$ and $M_{u} = \bar\beta\tilde{Y}$ in the frame of \claimref{synthesis.rung-field-coupling-to-its-conjugate-is-the-second-family-matrix}, which gives the determinants. From $M_{d} = \beta Y + \bar\alpha\tilde{Y}$ and $M_{u} = \alpha Y + \bar\beta\tilde{Y}$, $\beta M_{u} - \alpha M_{d} = (|\beta|^{2} - |\alpha|^{2})\tilde{Y}$ and $\bar\beta M_{d} - \bar\alpha M_{u} = (|\beta|^{2} - |\alpha|^{2})Y$. If $\tilde{Y}x = 0$ for a unit vector $x$ then $\beta M_{u}x = \alpha M_{d}x$, so $|\beta|$ times the least singular value of $M_{u}$ is at most $|\alpha|$ times the greatest of $M_{d}$; likewise for $Y$ with the roles exchanged. The check computes the whole solution space of the conditions as an exact kernel for three invertible pairs ($n = 2, 3$; symmetric and not) and finds $a_{3} = 0$ and $q_{1} = q_{2}$ on it.

\emph{(v).} For a group element with family rotations $U_{Q}$, $U_{u}$, $U_{d}$ and $g = e^{tK_{\Phi}}$ the invariance reads $U_{Q}^{\mathsf T}M_{u}(g\Phi)U_{u} = M_{u}(\Phi)$ and the same for $M_{d}$ with $U_{d}$. With $P(\Phi) = \det M_{u}(\Phi)\det M_{d}(\Phi) = p(\Phi_{1}, \Phi_{2})\,p(\bar\Phi_{2}, \bar\Phi_{1})$, taking determinants gives $P(e^{tK_{\Phi}}\Phi) = e^{-tS}P(\Phi)$, $S = I_{4}a_{3}$. Let $a_{3} \neq 0$ and $p \not\equiv 0$. Choose unitary coordinates $\zeta$ in which $K_{\Phi}$ is diagonal with entries $I_{4}q'_{1}$, $I_{4}q'_{2}$. Then $P = \hat{p}(\zeta)\,\hat{r}(\bar\zeta)$ with $\hat{p}$, $\hat{r}$ binary forms of degree $n$, and the monomial $\zeta_{1}^{j}\zeta_{2}^{n-j}\bar\zeta_{1}^{i}\bar\zeta_{2}^{n-i}$ has weight $I_{4}(j - i)(q'_{1} - q'_{2})$. Since $P$ has the single weight $-I_{4}a_{3} \neq 0$, $q'_{1} \neq q'_{2}$ and $j - i$ is the same for every pair of monomials of $\hat{p}$ and $\hat{r}$: each is a single monomial, $\hat{p} = c\,\zeta_{1}^{j}\zeta_{2}^{n-j}$, $\hat{r} = c'\bar\zeta_{1}^{i}\bar\zeta_{2}^{n-i}$, $i \neq j$. So $p$ vanishes exactly on the two orthogonal lines $\ell_{1} = \{\zeta_{1} = 0\}$, $\ell_{2} = \{\zeta_{2} = 0\}$, to orders $j$ and $n - j$. Now $\hat{r}$ is $\Phi \mapsto p(s\Phi)$ with $s(\Phi_{1}, \Phi_{2}) = (\bar\Phi_{2}, \bar\Phi_{1})$, an antilinear norm-preserving involution, so it vanishes on $s\ell_{1}$ to order $j$ and on $s\ell_{2}$ to order $n - j$, and by its monomial form on $\ell_{1}$ to order $i$ and on $\ell_{2}$ to order $n - i$. As $i \neq j$, $s\ell_{1} \neq \ell_{1}$ where $j > 0$ and $s\ell_{2} \neq \ell_{2}$ where $j < n$, so $s$ carries one of the two lines to the other, which is its orthogonal complement, and $i = n - j$. For a line spanned by $(a, b)$ the Hermitian product with $s(a, b) = (\bar{b}, \bar{a})$ is $2\bar{a}\bar{b}$, which vanishes only if $a = 0$ or $b = 0$: the lines are the coordinate lines of $\Phi$. Hence $K_{\Phi}$ is diagonal in $(\Phi_{1}, \Phi_{2})$, $p(a, b) = c\,a^{j}b^{n-j}$ up to exchanging $j$ and $n - j$ with the two weights, $2j \neq n$, and $-a_{3} = (j - i)(q_{1} - q_{2}) = (2j - n)(q_{1} - q_{2})$. Then $|\det M_{u}\det M_{d}| = |p(\alpha, \bar\beta)\,p(\beta, \bar\alpha)| = |c|^{2}|\alpha\beta|^{n}$, which is $|c|^{2}H^{2n}$ where $\alpha\beta = -H^{2}$ (\claimref{synthesis.rung-field-on-the-lift-locus}) and is $0$ at $\alpha = 0$ because $n \geq 1$. The three counter-models are computed in the check: in each the solution space contains a generator with $a_{3} = -(2j - n)(q_{1} - q_{2}) \neq 0$, and $p$ is the stated monomial ($j = n$ for $\tilde{Y} = 0$ and for $Y = 1$ with $\tilde{Y}$ nilpotent, since $\det(a + b\tilde{Y}) = a^{n}$).

\emph{(vi).} $B(e_{S}, e_{U}) \neq 0$ only for $U$ the complement of $S$ (\claimref{standard-model.invariant-pairing-five-modes}). A term $B(\psi, a\varphi)\,z$, with $a$ a slot operator and $z$ the field coefficient it carries, is conserved when the charges of the two states and of $z$ add to zero. In $B(\psi, c\varphi)$ the creation operators carry $\alpha$ and the contractions $\beta$: these are the four rows of clause (ii) of \claimref{synthesis.abelian-charges-of-the-sixteen-and-the-terms-that-break-them}, with a four-dimensional kernel on which $A_{3} = -2(h_{v} + h_{w})$. In $B(\psi, c^{*}\varphi)$, $c^{*} = \bar\beta\cdot v + \bar\alpha\cdot w$, the creation operators carry $\bar\beta$, of charge $-h_{w}$, and the contractions $\bar\alpha$, of charge $-h_{v}$: the same species sums with $-h_{w}$ in place of $h_{v}$ and $-h_{v}$ in place of $h_{w}$. Subtracting a row of one coupling from the row of the other on the same species gives $h_{v} + h_{w} = 0$, the hypothesis of the cited clause, so the kernel is its $\mathrm{span}\{Y, n_{B}, n_{L}\}$, on which $A_{3} = 0$ by its clause (v). The unit with $h_{v} = h_{w} = -2$ has total charge $1 + 1 - 2 = 0$ on the terms of the first coupling and $1 + 1 + 2 = 4$ on those of the second, so $\exp(\phi\,\cdot)$ of it keeps the second exactly when $4\phi$ is a whole number of turns. At a quarter turn it multiplies each of the sixteen by $e^{\pm I_{4}\pi/2} = \pm I_{4}$ and the slot coefficients by $e^{\mp I_{4}\pi} = -1$. The pseudoscalar $\omega$ acts on the sixteen by $-I$, with $I$ the unit of the module, and anticommutes with the generators, hence acts by $-1$ on the slot vectors (\claimref{synthesis.order-four-elements-with-square-the-grading-and-the-majorana-element}(ii), (iii)): it is the inverse quarter turn. $A_{3} = 8$ for the unit is clause (ii) of the cited charge list. The check recomputes the kernels and the charges.

\emph{(vii).} Under the axial phase $|\alpha|^{2}$ and $|\beta|^{2}$ are fixed, so $\mathsf{N}$ and $m$ are, and $Q = \alpha^{\mathsf T}\beta$ is multiplied by $e^{2I_{4}\phi}$, so $|Q|$ is fixed. By clause (1) of \claimref{synthesis.rung-field-d-term-action-and-tree-masses-on-the-lift-locus}, $\mu_{Y} = \tfrac12 m$ and $\sum_{j}\mu_{j}^{2} = \tfrac14(\mathsf{N}^{2} - |Q|^{2})$, so $V_{D}$ is fixed, with $W(\mathsf{N})$ and $k_{10}|Q|^{2}$; the terms $|\alpha|^{2}\operatorname{Re}Q$, $\operatorname{Re}Q^{2}$ and the others with $k_{4}, \dots, k_{9}$ of \claimref{synthesis.linear-rung-field-has-a-neutral-normal-pair-in-the-scalon-sum}(3) are not, and neither is the equation $Q = -H^{2}$ with $H \neq 0$. Hypercharge multiplies $\alpha$ by $e^{I_{4}\phi}$ and $\beta$ by $e^{-I_{4}\phi}$. The axial phase by $\phi$ followed by hypercharge by $\phi$ fixes $\beta$ and multiplies $\alpha$ by $e^{2I_{4}\phi}$, so it fixes the point $\alpha = 0$, $\beta = (-\lambda, 0)$ and acts on the tangent directions of $\alpha_{1}$, the normal pair of clause (2) there, as a rotation. A potential invariant under it has a Hessian at that point commuting with the rotation, so the two directions have equal masses: with $k_{8} = k_{9} = 0$ this is $r = 0$ in the stored formula $2W' + 4\mathsf{N}^{*}(k_{10} \pm r)$. The stored spectrum there has three zero modes, the gauge orbit, and no other.

\emph{(viii).} At $\alpha = 0$, in the frame with $\alpha$, $\beta$ along rung $4$, $\det M_{u} = p(\alpha_{1}, \bar\beta_{1})$ and $\det M_{d} = p(\beta_{1}, \bar\alpha_{1})$. For invertible $\tilde{Y}$, $\det(a Y + b\tilde{Y}) = b^{n}\det\tilde{Y}\,\det(1 + (a/b)\tilde{Y}^{-1}Y) = b^{n}\det\tilde{Y}\,(1 + (a/b)\tr(\tilde{Y}^{-1}Y) + \dots)$, so $\delta\log\det M_{u} = (\alpha_{1}/\bar\beta_{1})\tr(\tilde{Y}^{-1}Y)$; the same expansion in $b/a$ gives $\delta\log\det M_{d} = (\bar\alpha_{1}/\beta_{1})\tr(Y^{-1}\tilde{Y})$. The second term is the conjugate of $(\alpha_{1}/\bar\beta_{1})\,\overline{\tr(Y^{-1}\tilde{Y})}$, so the $I_{4}$ part of the sum is that of $(\alpha_{1}/\bar\beta_{1})\bigl(\tr(\tilde{Y}^{-1}Y) - \overline{\tr(Y^{-1}\tilde{Y})}\bigr)$, which vanishes for every $\alpha_{1}$ only if the bracket does. The check verifies the two coefficients on an exact pair.

\emph{The Reading} uses (i) to (viii) as stated; the two literature statements it names, the susceptibility and the review of axion searches are not stored Facts. For the estimate in (c): a field of mass $M$ coupled to the angle with gradient $\partial\bar\theta/\partial\alpha_{1}$ of size $|t|/\lambda$, $t$ the bracket of (viii), in a vacuum energy $\tfrac12\chi\bar\theta^{2}$ is displaced so as to change $\bar\theta$ by the relative part $\chi|t|^{2}/(\lambda^{2}M^{2})$ up to the normalisation of its kinetic term; with $\chi \approx (0.075\ \mathrm{GeV})^{4}$, $\lambda \approx 246$ GeV, $M \approx 3000$ GeV and $|t|$ of order the ratio of the top to the bottom mass, about $50$, this is about $3 \times 10^{-5} \cdot 2500/(6 \times 10^{4} \cdot 9 \times 10^{6}) \approx 10^{-13}$.
